%% The first command in your LaTeX source must be the \documentclass command.
\documentclass[
% twocolumn,
% hf,
]{ceurart}

\sloppy

\usepackage{booktabs}
\usepackage{tabularx}
\usepackage{array}
\usepackage{amsmath}
\usepackage{listings}
\lstset{breaklines=true,basicstyle=\ttfamily\small}

\newcolumntype{Y}{>{\raggedright\arraybackslash}X}
\newcommand{\code}[1]{\texttt{\detokenize{#1}}}

\begin{document}

\copyrightyear{2026}
\copyrightclause{Copyright for this paper by its authors.
  Use permitted under Creative Commons License Attribution 4.0
  International (CC BY 4.0).}

\conference{SEMANTiCS 2026 Posters and Demos Track}

\title{SEMANTiCS Scope Explorer: Query-Conditioned KG-RAG Summaries of Czech Public-Health Tables}

% CEUR and SEMANTiCS single-anonymous review require natural-person authors
% with meaningful affiliations. Replace this draft author block before
% submission; organizations or teams cannot be authors in CEUR-WS.
\author[1]{SEMANTiCS Project Team}
\address[1]{SEMANTiCS project workspace}

\begin{abstract}
SEMANTiCS Scope Explorer is a reproducible demo for turning Czech natural
language questions over large public-health tables into inspectable aggregate
SQL summaries. To our knowledge, it is the first RAG/KG-RAG analysis demo built
for this exact combination of Czech NZIS infectious-disease aggregates and
MKN-10-CZ code-list data. The system compares a flat RAG baseline with KG-RAG
over the same 272,858-row DuckDB fact table, so the difference is scope
construction rather than database access. Flat RAG selects table-present
values from retrieved text. KG-RAG adds typed MKN-10-CZ hierarchy expansion
before SQL, converting disease-family phrases into complete diagnosis-code
scopes. On ten saved live demo queries, KG-RAG
improves diagnosis micro-F1 from 0.545 to 0.935, macro-F1 from 0.567 to
0.921, exact diagnosis scopes from 1/10 to 8/10, and complete-filter accuracy
from 1/10 to 6/10, while producing no false-positive diagnosis codes. The demo
exposes the retrieved evidence, graph expansion, generated SQL, and returned
table, making the semantic gain and remaining failure modes directly auditable.
\end{abstract}

\begin{keywords}
retrieval-augmented generation \sep
knowledge graphs \sep
KG-RAG \sep
semantic retrieval \sep
text-to-SQL \sep
public-health open data
\sep tabular summarization
\end{keywords}

\maketitle

\section{Problem and Contribution}

Public-health table exploration is often framed as document search, but the
hard part is usually scope construction. Even a simple user phrase such as
``děti'' (children) is not a database predicate. The source table stores
diagnosis codes, official age-group and region labels, months, years, and one
reported aggregate measure. A fluent answer is not useful if it silently
omits sibling diagnoses or applies informal values that cannot be executed.

SEMANTiCS Scope Explorer treats question answering as query-conditioned
summarization over a table, not as prose generation over retrieved chunks. The
demo contributes a shared DuckDB answer layer, a flat RAG baseline, a KG-RAG
variant with typed MKN-10-CZ hierarchy expansion, and a saved ten-query
evaluation. To our knowledge, no prior demo has reported RAG-versus-KG-RAG
scope accuracy on this exact NZIS infectious-disease source joined to the
Czech MKN-10-CZ terminology. The comparison is controlled: both methods query
the same \code{fact_infectious_disease_cases_enriched} table, and both return
selected entities, generated SQL, and a rendered result table. The reported
measure is always \code{reported_case_count} (source \code{pocet_pripadu});
it is not interpreted as unique patients, incidence, prevalence, causal
evidence, or clinical guidance.

\section{Architecture}

The current SQL answer layer intentionally focuses on a newly built
infectious-disease fact table rather than claiming to answer over every
profiled source. The final answer table contains 272,858 infectious-disease
rows from 2018--2025, enriched with MKN-10-CZ terminology, derived MKN
blocks, chapter labels, age intervals, region labels, sex labels, and a
stable measure. MKN-10-CZ (39,136 rows) and the NZIS age-group code list (228
rows) are joined into this table. The 1,418,107-row hospitalization dataset is
loaded into DuckDB and represented in the KG through diagnosis, year, and
measure context, but it is not yet joined into the final SQL answer table.
This keeps the RAG/KG-RAG comparison controlled on the new infectious-disease
dataset build; cross-fact analyses with hospitalization are the next step.
All sources come from Czech public-health open data \cite{nzip_sources}.
DuckDB provides the local analytical execution layer
\cite{raasveldt2019duckdb}.

The AI Scope Builder has four shared runtime stages. First, it extracts
candidate spans from the user's Czech question. Second, it embeds each span
and compares it with a flat catalog of public values and grouping columns. The
database candidate score is
\[
\mathrm{score}_{DB}=0.86\,\cos(e_s,e_c)+0.14\,\mathrm{lex}(s,c),
\]
with candidates below cosine 0.20 discarded. Third, a model selects filters,
groupings, and context entities from candidate IDs. Deterministic safeguards
promote at most five diagnosis filters per disease/code span, require score
at least 0.28, and allow no more than a 0.14 drop from the best diagnosis
candidate. Fourth, SQL generation is constrained to a single read-only
\code{SELECT} over the shared final table. The AI methods are not allowed to
filter directly on MKN block or chapter columns; hierarchy must be converted
into explicit \code{diagnosis_code IN (...)} values.

Flat RAG uses 114 embedded documents, one per table-present infectious
diagnosis. KG-RAG uses 40,906 node documents and a graph with 40,906 nodes and
87,364 edges. The graph contains MKN chapters, derived MKN blocks, MKN codes,
infectious-disease diagnosis values, hospitalization diagnosis values, years,
months, measures, datasets, and columns. Its main expansion edges are
\code{CONTAINS_MKN_CODE}, \code{HAS_TABLE_DIAGNOSIS_VALUE}, and
\code{SAME_AS_MKN_CODE}; it also stores 2,418 weighted co-occurrence edges and
1,169 diagnosis profile-similarity edges.

KG-RAG inserts one extra step between flat entity selection and SQL. Selected
diagnosis seeds are grouped by MKN block and ranked by
\[
\mathrm{score}_{block}=0.55\,\frac{|C_b|}{|C|}
+0.25\,\frac{\sum_{d\in C_b}w_d}{\sum_{d\in C}w_d}
+0.20\,\mathrm{lex}(q,b).
\]
Here $C$ is the selected seed set, $C_b$ its subset inside block $b$, and
$w_d$ is the reported-case total for diagnosis $d$. If block lexical
similarity is at least 0.65, or no seed code is already present in the table
for that block, the whole table-present block is used; otherwise the system
uses a contiguous seed range or seed-only expansion. These ratios are
deterministic demo parameters chosen through interactive user-experience
tuning with the live demo. They are not learned model parameters and were not
fit by optimizing the evaluation metrics.

\section{Evaluation}

The saved evaluation in \code{dg_evaluation} replays ten Czech user queries
against live API outputs for RAG and KG-RAG. Gold diagnosis groups such as
A00--A09, B00--B09, B15--B19, and J12--J18 are resolved by intersecting the
official range with diagnosis codes present in the DuckDB fact table. Other
gold filters are official year, age-group, and region values from the
evaluation JSON and DuckDB lookups. The scorer parses SQL \code{WHERE}
clauses first; if an expected filter is not explicit in SQL but appears in the
returned rows, the distinct returned row values are used. Complete-filter
accuracy is binary per question: every expected filter column must have
exactly the expected values and no unexpected filter column.

\begin{table}[h]
\caption{Diagnosis-scope effectiveness. Values are precision/recall/F1; Ex is exact diagnosis-set match.}
\label{tab:diagnosis-results}
\scriptsize
\begin{tabularx}{\linewidth}{llcccc}
\toprule
ID & Target & RAG & Ex & KG-RAG & Ex \\
\midrule
EQ-01 & A00--A09 & .60/.33/.43 & 0 & 1.00/1.00/1.00 & 1 \\
EQ-02 & A00--A09 & .60/.33/.43 & 0 & 1.00/1.00/1.00 & 1 \\
EQ-03 & A00--A09 & .60/.33/.43 & 0 & 1.00/1.00/1.00 & 1 \\
EQ-04 & A00--A09 & .60/.33/.43 & 0 & 1.00/1.00/1.00 & 1 \\
EQ-05 & A00--A09 & .60/.33/.43 & 0 & 1.00/1.00/1.00 & 1 \\
EQ-06 & A00--A09 & .60/.33/.43 & 0 & 1.00/.56/.71 & 0 \\
EQ-07 & B00--B09 & .60/.33/.43 & 0 & 1.00/.33/.50 & 0 \\
EQ-08 & B15--B19 & 1.00/1.00/1.00 & 1 & 1.00/1.00/1.00 & 1 \\
EQ-09 & J12--J18 & 1.00/.71/.83 & 0 & 1.00/1.00/1.00 & 1 \\
EQ-10 & J12--J18 & 1.00/.71/.83 & 0 & 1.00/1.00/1.00 & 1 \\
\midrule
All & macro & .72/.48/.57 & 1/10 & 1.00/.89/.92 & 8/10 \\
All & micro & .72/.44/.55 & -- & 1.00/.88/.94 & -- \\
\bottomrule
\end{tabularx}
\end{table}

Table~\ref{tab:diagnosis-results} shows the central result. KG-RAG doubles
micro recall from 0.439 to 0.878 and raises exact diagnosis scopes from one
query to eight. More importantly for a public demo, KG-RAG produces zero
false-positive diagnosis codes across the ten queries, while flat RAG adds 14
wrong diagnosis codes. The main win is not a nicer answer sentence; it is the
correct set of values over which a huge table is summarized.

\begin{table}[h]
\caption{Complete-filter effectiveness. A correct run must match all expected diagnosis, time, age, and region filters.}
\label{tab:filter-results}
\scriptsize
\begin{tabularx}{\linewidth}{llccY}
\toprule
ID & Expected filters & RAG & KG-RAG & KG-RAG residual when incorrect \\
\midrule
EQ-01 & diagnosis & 0 & 1 & -- \\
EQ-02 & diagnosis + year & 0 & 1 & -- \\
EQ-03 & diagnosis + age filter & 0 & 0 & age filter broadened to all groups \\
EQ-04 & diagnosis + age filter & 0 & 1 & -- \\
EQ-05 & diagnosis + region filter & 0 & 0 & region filter broadened to all regions \\
EQ-06 & diagnosis + region filter & 0 & 0 & region kept; diagnosis scope partial \\
EQ-07 & diagnosis & 0 & 0 & B00--B09 scope partial \\
EQ-08 & diagnosis & 1 & 1 & -- \\
EQ-09 & diagnosis & 0 & 1 & -- \\
EQ-10 & diagnosis & 0 & 1 & -- \\
\midrule
All & complete filters & 1/10 & 6/10 & -- \\
\bottomrule
\end{tabularx}
\end{table}

Table~\ref{tab:filter-results} gives the stricter table-summarization view:
all filters must be right, not only disease family. KG-RAG still improves
accuracy from 0.10 to 0.60. The remaining failures are informative. EQ-03 and
EQ-05 are not diagnosis-hierarchy failures; they expose that age and regional
aliases are currently handled by the shared flat scope builder rather than by
a graph of age and region concepts. EQ-06 and EQ-07 are partial hierarchy
expansions caused by the selected seeds and conservative block-expansion
policy. The demo therefore shows both a strong effectiveness gain and an
auditable path for the next engineering iteration.

\section{Discussion and Conclusion}

The results support the paper's core claim: KG-RAG is valuable for
query-conditioned table summarization when a user phrase denotes a coded
family rather than a literal table value. The A00--A09 and J12--J18 cases show
why graph closure matters. Many flat RAG approaches retrieve plausible
high-precision seeds but fail to retrieve the complete controlled-vocabulary
set, so recall becomes the weak point even when the top retrieved labels look
correct. Our flat baseline follows this pattern: it often finds relevant
diagnoses, but it misses sibling codes needed for complete table
summarization and sometimes introduces outsiders. KG-RAG turns the same
initial evidence into an explicit code set, then lets DuckDB compute the
answer. This separation of semantic scope construction from SQL execution is
what makes the demo inspectable and repeatable.

The current limitations are also clear. MKN blocks are configured from known
ranges because the source CSV provides chapters but no explicit block rows.
The graph does not yet model age concepts or regional aliases; those filters
use shared flat candidate selection. The SQL answer layer also does not yet
expose all loaded datasets as answer facts: hospitalization is available as
source and graph context, while infectious-disease reported cases are the
current summarized measure. The numeric ratios in candidate ranking and
hierarchy ranking are transparent demo-tuned settings, not trained
parameters. Future work should add typed age/region concept nodes, extend the
new dataset build into multi-fact hospitalization analyses, learn or calibrate
the ratios on held-out questions, and evaluate larger multilingual query
sets. The system uses aggregate public-health data and does not provide
clinical advice, causal claims, or patient-level inference.

\section*{Declaration on Generative AI}
During preparation of this manuscript draft, the authors used an AI coding
assistant to transform project artifacts into \LaTeX{} prose and to check
compilation. The authors are responsible for reviewing, editing, and
validating the final content before submission.

\bibliography{semantics-paper}

\end{document}
